\subsection{The result and its scope}
This construction compiles an explicitly given nonempty quadratic system over
\(\mathbb Z\) into a nonnegative polynomial of degree at most four and
coefficient height one. Source rows are collected and nonzero, and every
source coordinate occurs; the supplied systems satisfy these conditions. Its
new step is to contract aliases using the \emph{actual incidence energy},
rather than protecting every possible pair of factor monomials. Two further
improvements factor affine squares before coefficient normalization and
replace certain three-coordinate equality cycles by a two-coordinate
quartic. The supplied coefficient files, not an unexpanded circuit alone,
define the final polynomials.

The correspondence proved below is a bijection between integer zero sets.
It is unconditional finite algebra. Applying it to the supplied calibrated
source establishes the same ZFC-provable equivalence between unsolvability
and \(\operatorname{Con}(\mathrm{ZFC})\) as that source. Neither the
construction nor its finite checks prove either consistency or universal
nonvanishing. No global optimum, literature-wide priority, or new
end-to-end formal-kernel consistency calibration is claimed.

\subsection{Exact affine refactorization}
Let the original system be \(R_r(x)=0\), where each \(R_r\) is integral
of total degree at most two. Suppose a homogeneous quadratic part is
\(\epsilon L(x)^2\), where \(\epsilon\in\{-1,1\}\) and \(L\) is
an integral linear form. Choose an integer \(k\), introduce \(z\), and
replace that row by
\[
 z-L(x)-k=0,\qquad
 \epsilon z^2+R_r(x)-\epsilon(L(x)+k)^2=0.
\]
The second displayed expression is quadratic because the homogeneous
quadratic parts cancel. A previously introduced coordinate for the same
affine form may be reused. Every added definition is monic in a fresh
coordinate and depends only on original coordinates. Consequently, for
every original tuple there is exactly one extension, namely
\(z=L(x)+k\). Substitution proves the exact pullback identity for every
replaced row. Reading the original coordinates is its inverse on zeros.
No division, sign assumption, inequality, or domain change is used.

The producer chooses candidate forms by exact integer square tests and
coefficient comparison, never by numerical factorization. The independent
checker instead reconstructs every affine definition and proves every
original-row pullback coefficient by coefficient. Correctness therefore
does not depend on the heuristic selecting the sparsest shift.

\subsection{Unit coefficient normalization}
Write the refactored source coordinates as \(X\). For every needed positive
coefficient \(c>1\), and recursively for \(\lfloor c/2\rfloor>1\),
introduce a coordinate \(t_c\) and the monic equation
\[
 t_c-t_{\lfloor c/2\rfloor}-t_{\lfloor c/2\rfloor}
       -(c\bmod 2)=0,
 \qquad t_1:=1.
\]
The two equal summands are retained as separate signed occurrences. In
increasing order of \(c\), these equations force \(t_c=c\). For each
needed product \(cX_iX_j\), introduce a shared scale coordinate
\(s_{c,i}\) with \(s_{c,i}-t_cX_i=0\). Replace a signed source term of
absolute coefficient \(c\) by, respectively,
\[
  t_c,\qquad t_cX_i,\qquad s_{c,i}X_j
\]
for constant, linear, and quadratic terms. Terms of absolute coefficient
one need no change. All resulting rows have the form
\[
 S_r(y)=\sum_{j=0}^{s_r-1}\epsilon_{rj}m_{rj}(y),
 \qquad \epsilon_{rj}\in\{-1,1\},\quad \deg m_{rj}\leq2.
\]
The new coordinates are uniquely determined by the monic constructors.
The normalization checker verifies the topological pivots, their complete
coverage, and the exact polynomial pullback of every source row. Thus this
step is also a bijection of integer zero sets.

\subsection{A shared incidence graph}
Fix a positive integer \(H\). Make one hub block of \(H\) vertices. For
row \(r\), make a row block of \(b_r=\lceil s_r/H\rceil\) vertices.
Connect the vertices of each block by a tree with a different unrestricted
integer flow coordinate on each internal edge. The number of free flow
coordinates is exactly
\[
 F(H)=H-1+\sum_r(b_r-1).
\]
For each signed occurrence, add an edge directed from its row block to
the hub block, with prescribed value \(\epsilon_{rj}M_{rj}\).
The explicit placement in the certificate makes these edges distinct;
there are no loops or parallel edges. Each \(M_{rj}\) is a homogeneous
quadratic monomial in factor aliases. A quadratic source term uses aliases
of its two logical factors. A linear term uses an alias of its logical
factor times a unit alias. A constant term uses two unit aliases. The
factors of a square may be the same coordinate.

For each vertex \(v\), let \(d_v\) be the outgoing minus incoming edge
sum. Define
\[
 E=\frac12\sum_v d_v^2.
\]
Before alias contraction, distinct formal edge values have diagonal
coefficient one, since every edge has two endpoints. Distinct edges meet
in at most one vertex, so their cross coefficient is zero or \(\pm1\).
After monomial substitution, coefficients may collide. This is a
\emph{coefficient} issue, not a zero-set issue: the exact sum-of-squares
identity remains valid under every substitution.

\subsection{Incidence local contraction}
Every factor coordinate has a logical value: one normalized coordinate
\(y_i\), or the distinguished unit value. Contract only coordinates with
the same logical value. The producer tries a substitution \(b\mapsto a\)
and collects \emph{all} affected terms of \(E\). It accepts the substitution
only when each resulting coefficient belongs to \(\{-1,0,1\}\).
Opposite signs may cancel; zero coefficients are deleted. No C4-free
condition is imposed on the factor graph, and square aliases may merge
with other aliases when the exact coefficient test permits it.

This is strictly less restrictive than requiring uniqueness of the
product of \emph{every} pair of source monomials. Only a source-edge
square or a pair of incident edges contributes to \(E\). A collision
between two products not simultaneously present in that support is
irrelevant. Even a collision in the support can be harmless if its signed
coefficients cancel. The independent checker does not trust the contraction
history: it reconstructs the final vertex divergences, expands
\(\sum_vd_v^2\), and verifies the entire collected polynomial.

Contraction does not change either direction of the intended zero-set
map. At every intended zero, aliases with the same logical value were
already equal. Conversely, the following equality penalties force all
remaining aliases in each logical class to agree. Every factor occurrence
is independently checked against its normalized source monomial.

\subsection{Equality and unit penalties}
For a logical class of \(k\geq3\) distinct aliases, use the cycle energy
\[
 C(a_1,\ldots,a_k)=\sum_{i=1}^k a_i^2-
                  \sum_{i=1}^k a_i a_{i+1},\qquad a_{k+1}=a_1.
\]
Twice this energy is \(\sum_i(a_i-a_{i+1})^2\). Its coefficients have
height one, and it vanishes precisely when all aliases agree. A singleton
requires no penalty. A two-member class can be completed to a three-cycle
with one new dummy alias.

There is a better option for a suitable two-member class:
\begin{align*}
 Q(a,b)&=a^4+b^4-a^3b-ab^3\\
       &=(a-b)^2(a^2+ab+b^2).
\end{align*}
Since \(a^2+ab+b^2\) is positive except at \((0,0)\),
\(Q\geq0\) and \(Q=0\) if and only if \(a=b\). An exact
sum-of-squares identity, also used by the checker, is
\[
 2Q=(a^2-b^2)^2+(a^2-ab)^2+(ab-b^2)^2.
\]
If neither \(a^2\) nor \(b^2\) occurs as a prescribed edge monomial,
none of the four monomials of \(Q\) can occur in \(E\): a product
producing \(a^4\) needs an \(a^2\) factor, and one producing
\(a^3b\) needs an \(a^2\) factor; the other cases are symmetric.
The lower-degree flow terms cannot collide either. Distinct logical
classes are disjoint. Hence this criterion permits replacing the dummy
and its cycle by \(Q\), saving one variable and two monomials while
preserving height one. The emitted polynomial is checked in full even
when this structural sufficient condition is available.

Let \(u_1,\ldots,u_\ell\), with \(\ell\geq2\), be all unit aliases,
adding dummies if needed. Use
\[
 U=1+\sum_{i=1}^{\ell}u_i^2-u_1-u_\ell-
                      \sum_{i=1}^{\ell-1}u_iu_{i+1}.
\]
Twice \(U\) equals
\((u_1-1)^2+\sum_{i<\ell}(u_i-u_{i+1})^2+(u_\ell-1)^2\).
Thus \(U\geq0\), with equality precisely when every unit alias is one.
Unit, ordinary-alias, and free-flow namespaces are disjoint.

The final polynomial is
\[
 \boxed{\displaystyle P=E+\sum_{\text{cycle classes}}C+
                         \sum_{\text{pair classes}}Q+U.}
\]
It is nonnegative over \(\mathbb R\). The source-edge products have
degree two, so \(\deg P\leq4\). The only constant term is the one
in \(U\). Full collected expansion proves that every nonzero coefficient
is \(\pm1\), that the exact degree is four for each delivered endpoint,
and that every counted variable actually occurs.

\subsection{Explicit forward witness map}
Given original integer coordinates \(x\), first evaluate all affine
refactor definitions, then all constant and scale definitions. Assign
every ordinary alias and its dummy the corresponding normalized value,
and every unit alias or dummy the value one. All equality and unit
penalties then vanish identically.

Root every internal block tree at its first listed vertex. Let \(b_v\)
be the divergence contributed at \(v\) by prescribed edges alone.
For an internal tree edge joining a child \(v\) to its parent, let
\(\sigma_v\in\{-1,1\}\) be that edge's incidence sign at the child.
Set its flow to
\[
 w_v=-\sigma_v\sum_{t\text{ in the subtree rooted at }v}b_t.
\]
This is an explicit integer polynomial in the source coordinates. Summing
subtree divergences shows that every nonroot divergence is zero. The
root divergence is the sum of prescribed divergences over its block.
For row block \(r\) it is \(S_r\); for the hub block it is
\(-\sum_rS_r\). Constructor rows pull back to zero, while all other
rows pull back to the original residuals. Therefore the full lift \(L\)
satisfies the universal identity
\[
 P(L(x))=\frac12\left(\sum_rR_r(x)^2+
                                \left(\sum_rR_r(x)\right)^2\right).
\]
In particular an original zero gives a final integer zero. The displayed
identity is \emph{not} generally \(P(L(x))=\sum_rR_r(x)^2\).
Its integrality follows either from the integer polynomial construction
or from \(n^2\equiv n\pmod2\).

\subsection{Explicit inverse and uniqueness}
Suppose \(P(z)=0\) for an integer tuple \(z\). Every summand in the
boxed definition is nonnegative, so each is zero. The equality penalties
identify all aliases in each logical class, and the unit path forces all
unit values to one. Read one listed alias representative from each
logical class to obtain \(y\). Every vertex divergence vanishes because
\(E=0\). Summing divergences over row block \(r\) cancels its internal
edges and yields exactly \(S_r(y)=0\). Topological normalization forces
each constant and scale coordinate to its prescribed value; affine
refactor definitions force each added affine coordinate. Reading the
original coordinates therefore gives \(R_r(x)=0\) for every \(r\).

All remaining coordinates are uniquely determined by this \(x\). The
constructor coordinates, aliases, units, and dummies are unique by their
definitions. On each internal tree, eliminating a leaf uniquely determines
its incident free flow, since its coefficient is \(\pm1\). Continuing
to the root gives precisely the subtree formula above. Thus reading the
original coordinates is inverse to \(L\) on zero sets, proving a
bijection, not just equisolvability.

\subsection{Reproducibility and evidence boundaries}
The literal source equations, affine certificates, graph certificates,
and fully collected polynomial files determine every mathematical object.
The producer uses exact standard-library integer arithmetic. The independent
checkers import no producer, reconstruct all affine source pullbacks,
verify every alias occurrence, check the simple incidence graph and every
internal tree, and expand twice the final polynomial as an actual sum of
squares. Both ordinary Python and \texttt{python -O} runs use explicit
exceptions rather than removable assertions. Deliberate corruption tests
exercise coefficient signs, source hashes, source-row coverage, logical
factor bindings, missing edges, alias coverage, unit endpoints, and metrics.
Finite off-zero evaluations are supplementary implementation tests only;
the preceding algebra establishes the unbounded witness maps.

The incoming shared-incidence construction, coefficient normalization,
cycle and unit energies, and their supplied source calibration retain
their original attribution. The contribution here is incidence-local
coefficient-aware contraction, the collision-safe quartic pair penalty,
and their composition with exact affine refactorization and the new source.
Comparisons are tied to the inspected archives and exact source hashes.
